\documentclass[
% papersize
a4paper,
%fond size
%12pt,
11pt,
%page print style
twoside,
%oneside,
openright]{book}

%for PSTRicks
\usepackage[pdf]{pstricks}
\usepackage[off]{auto-pst-pdf}

%graphics included tikz & pfg
\usepackage{graphicx}
\usepackage{epsfig,epic,eepic}
\usepackage{psfrag}
\usepackage[USenglish,german]{babel}
\usepackage{color,pstricks}
% To allow me to use letters with accent
\usepackage[utf8]{inputenc}

% Fonts
%\usepackage{lmodern}
%from Ruben \usepackage[T1]{fontenc}
\usepackage{times}
%\usepackage[scaled]{helvet}
%\usepackage[sc]{mathpazo} % math fonts for Palatino, see
                          % http://www.tug.dk/FontCatalogue/palatino/
%\usepackage{palatino}

% Math packages
\usepackage{amsmath,amssymb,amsfonts,amsthm}
\usepackage{url}
\usepackage[ruled,noline]{algorithm2e}
% More compact lists
\usepackage{paralist}
% Fullpage
\usepackage[in,headings]{fullpage}
% But then adapt spacing
\usepackage{setspace}
\setstretch{1.3}

%for nice tables
\usepackage{multirow}
%\usepackage{booktabs}
\usepackage{tikz}
\usepackage{ctable}

%for nice item lists
\usepackage{paralist}

%usepackage comment for block comments
\usepackage{comment}

%nice symbols to draw arrow (yes) and -(no)
\usepackage{pifont}
\newcommand{\yes}{\checkmark}
\newcommand{\no}{\textendash}

%sub-picture numbering
%\usepackage{subcaption}
\usepackage{subfigure}

\usepackage[right]{eurosym}


%annotate TODO stuff
\usepackage[colorinlistoftodos, textwidth=4cm, shadow]{todonotes}
\usepackage{lscape}

%pseudocode package
\usepackage{pseudocode}

%to insert some clickable links
\usepackage{hyperref}
\definecolor{darkblue}{rgb}{0,0,.5}
\definecolor{black}{rgb}{0,0,0}
\hypersetup{colorlinks=true, breaklinks=true, citecolor=black, linkcolor=black, menucolor=black, urlcolor=black}

% \setlength{\oddsidemargin}{-0.3cm}
% \setlength{\evensidemargin}{-0.3cm}
% \setlength{\textwidth}{16cm}
% \setlength{\topmargin}{-0.8cm}
% \setlength{\textheight}{23.0cm}

% use some common definitions
\newcommand{\pref}[1]{Part~\ref{#1}}
\newcommand{\cref}[1]{Chapter~\ref{#1}}
\newcommand{\sref}[1]{Section~\ref{#1}}
\newcommand{\aref}[1]{Appendix~\ref{#1}}
\newcommand{\fref}[1]{Figure~\ref{#1}}
\newcommand{\tref}[1]{Table~\ref{#1}}
\newcommand{\eref}[1]{(\ref{#1})}

% More convenient sometimes
\newcommand{\ie}{i.\@e.\@ }
\newcommand{\eg}{e.\@g.\@ }

\newcommand{\us}{\,$\mu$s\xspace}

%add empty page macro
\newcommand{\blankpage}{
\newpage
\thispagestyle{empty}
\mbox{}
\newpage
}

% Some math definitions

%\newcommand{\indic}{1\hspace{-2.5mm}{1}}
%\newcommand{\indic}{\boldsymbol{1}}
\newcommand{\indic}[1]{1_{\{#1\}}}
\newcommand{\erfc}{\mathrm{erfc}}
\newcommand{\abs}[1]{\left\lvert #1 \right\rvert}
\newcommand{\lp}[1]{\left( #1 \right)}
\newcommand{\lb}[1]{\left\lbrack #1 \right\rbrack}
\newcommand{\lc}[1]{\left\lbrace #1 \right\rbrace}
\newcommand{\erf}[1]{\mathrm{erf}\lp{#1}}
\newtheorem{definition}{Definition}
\newtheorem{lemma}{Lemma}
\newtheorem{theorem}{Theorem}

\newcommand{\calI}{\mathcal I}
\newcommand{\calN}{\mathcal N}
\newcommand{\calP}{\mathcal P}
\newcommand{\calX}{\mathcal X}
\newcommand{\calC}{\mathcal C}

% Paper specific notation
\def\su{^{(u)}}
\def\s0{^{(0)}}
%
\def\Xr{X^{(r)}}
\def\Xt{X^{(t)}}
%
\def\Ns{N_{\textrm{s}}}
\def\Na{N_{\textrm{a}}}
%
\def\tacq{t_{\textrm{acq}}}
\def\ttx{t_{\textrm{tx}}}
\def\tfail{t_{\textrm{fail}}}
\def\tdrop{t_{\textrm{drop}}}
%
\def\tprop{t_{\textrm{prop}}}
\def\tpreamble{t_{\textrm{preamble}}}
\def\tdata{t_{\textrm{DATA}}}
\def\tack{t_{\textrm{ACK}}}
\def\tsigidle{t_{\textrm{SIGIDLE}}}
\def\tmaxbackoff{t_{\textrm{max backoff}}}
\def\tsendtimer{t_{\textrm{send timer}}}
\def\tidletimer{t_{\textrm{idle timer}}}
\def\tbackoff{t_{\textrm{backoff}}}
%
\def\SD{S_{\textrm{D}}}
\def\SI{S_{\textrm{I}}}
%
\def\Pacq{p_{\textrm{acq}}}
\def\Pfail{p_{\textrm{fail}}}
\def\Pbusy{P_{\textrm{busy}}}
\def\Pmd{P_{\textrm{MD}}}
\def\Pfa{P_{\textrm{FACQ}}}
%
\def\td{t_{\textrm{D}}}
\def\ti{t_{\textrm{I}}}
\def\Np{N_{\textrm{p}}}


\newcommand{\mand} {\mathrm{\;  and \; }}
\newcommand{\mfor} {\mathrm{\;  for \;  }}
\def\Ints{\mathbb{Z}}
\def\P{\mathbb{P}}
\def\E{\mathbb{E}}
\def\be{\begin{equation}}
\def\ee{\end{equation}}
\def\ben{\[}
\def\een{\]}
\def\bearn{\begin{eqnarray*}}
\def\eearn{\end{eqnarray*}}
\def\bear{\begin{eqnarray}}
\def\eear{\end{eqnarray}}
\def\barr{\begin{array}}
\def\earr{\end{array}}
\def\bel{\be \barr{l}}% equation array adjusted left, numbered
\def\eel{\earr\ee}
\def\beln{\ben \barr{l}}% equation array adjusted left, no number
\def\eeln{\earr\een}\def\be{\begin{equation}}
\def\ee{\end{equation}}
\def\ben{\[}
\def\een{\]}
\def\bearn{\begin{eqnarray*}}
\def\eearn{\end{eqnarray*}}
\def\bear{\begin{eqnarray}}
\def\eear{\end{eqnarray}}
\def\barr{\begin{array}}
\def\earr{\end{array}}
\def\bel{\be \barr{l}}% equation array adjusted left, numbered
\def\eel{\earr\ee}
\def\beln{\ben \barr{l}}% equation array adjusted left, no number
\def\eeln{\earr\een}


% Style headings + headers and footers

% Modify header and footers
\usepackage{fancyhdr}
\pagestyle{fancy}
\fancyhf{}
\renewcommand{\chaptermark}[1]{
  \markboth{\thechapter.\ #1}{}
%  \markboth{#1}{}
}
\renewcommand{\sectionmark}[1]{
%  \markboth{\thechapter.\ #1}{}
  \markright{#1}
}

% \fancyhead[LE]{\normalfont\sffamily\thepage}
% \fancyhead[RE]{\normalfont\sffamily\leftmark}
% \fancyhead[LO]{\normalfont\sffamily\rightmark}
% \fancyhead[RO]{\normalfont\sffamily\thepage}
\fancyhead[LE]{\normalfont\thepage}
\fancyhead[RE]{\normalfont\leftmark}
\fancyhead[LO]{\normalfont\rightmark}
\fancyhead[RO]{\normalfont\thepage}

% Modify the fonts for the chapters, section, sub... and paragraphs
\usepackage{titlesec}

\titleformat{\chapter}[display]
%{\normalfont\Huge\sffamily}
{\normalfont\Huge}
{\filleft{\textbf{\LARGE{\chaptername} {\thechapter}}}}{1em}{\filleft\textbf}

\titleformat{\section}
{\normalfont\Large}
{\filright{\textbf{\thesection}}}{1em}{\filright\textbf}

\titleformat{\subsection}
{\normalfont\normalsize}
{\filright{\textbf{\thesubsection}}}{1.3em}{\filright\textbf}

\titleformat{\subsubsection}
{\normalfont\normalsize}
{}{0em}{\filright\textbf}

\titleformat{\paragraph}[runin]
{\normalfont\normalsize}
{}{0em}{\filright\textbf}

% Do it for theorems,... too
%\theoremstyle{definition}
%\newtheorem{define}{\textbf\normalfont\normalsize\sffamily{Definition}}[chapter]
%\theoremstyle{plain}
%\newtheorem{theorem}{\textbf\normalfont\normalsize\sffamily{Theorem}}[chapter]
%\newtheorem{lem}{\textbf\normalfont\normalsize\sffamily{Lemma}}[chapter]


% The larger it is, the less willing LaTeX is to split footnotes
\interfootnotelinepenalty=10000


% Removing master, slows down the compilation
\includeonly{
	title,
	abstract,
	acknowledgments,
	paper-collaborators,
	eid,
	chapter-1_introduction,
	chapter-2_related_work,
	chapter-3_measurement_tools,
	chapter-4_analyse_optimization,
	chapter-5_conclusion,
%	appendix,
}


% BEGINN of DOCUMENT
\begin{document}

\setlength{\headheight}{13.6pt}

%TITLE PAGE

% Titeseite

\begin{titlepage}
  \begin{center}
	  
    %\mbox{}

    %\vspace{1cm}

	\centering
	\includegraphics[width=3cm]{figures/tub.pdf}\\
	\vspace{0.8em}
	\LARGE 

	\textsc{Technische Universit\"at Berlin}

    \vspace{1cm}

    \sffamily \LARGE \textbf{My new Master thesitests titel comes here}

    \vspace{1.5cm}





    \normalsize vorgelegt von

    \vspace{.1cm}

    \large Alex the student (Dipl.-Ing.)

    \vspace{.8cm}

  %  von der Fakult�t IV - Elektrotechnik und Informatik der Technischen Universit�t Berlin


    \normalsize von der Fakult\"{a}t IV -- Elektrotechnik und Informatik\\
    %\vspace{.1cm}
    \normalsize der Technischen Universit\"{a}t Berlin\\
    %\vspace{.1cm}
    \normalsize zur Erlangung des akademischen Grades\\
    %\vspace{.1cm}
    \large \textsc{Diplom der Ingenieurwissenschaften (Dipl.-Ing.)}\\
    %\vspace{.1cm}
    \normalsize genemigte Diplomarbeit\\

    \vspace{1cm}

    \large \textbf{Diplomausschuss:}

    \vspace{.2cm}

		\begin{tabular}{rl}
		Pr\"ufer der Diplomarbeit:	& Prof. Dr.-Ing. XYZ, TU-Berlin\\
						& Dr.-Ing. XYZ, TU-Berlin\\
		\end{tabular}\\

    \vspace{2cm}
    \large Tag der wissenschaftlichen Verteidigung: 11.01.2013\\
    \vspace{0.5cm}
    \large Berlin 2013\\

  \end{center}
\end{titlepage}


\selectlanguage{USenglish}

%ABSTRACT_EN
\frontmatter \pagestyle{plain}
\setcounter{page}{1}
%\addcontentsline{toc}{chapter}{Resume}
\chapter*{Abstract}
\label{chap:abstract}

%
\textit Abstract .....




%ACKNOWLAGEMENTS
\addcontentsline{toc}{chapter}{Acknowledgments}
\chapter*{Acknowledgments}
\label{chap:ack}

%Acknowledgements for all the helping hands

\section*{}
At first, I want to thank ...





%PUBLICATIONS & COLLABORATORS
% \chapter*{Pre-Published Papers}
% \label{chap:prepublished_papers}

% Parts of this thesis are based on the following peer-reviewed papers that have already been published. All my collaborators are among my co-authors.
% %

% %%%%%%%%%%%%%%%%%%%%%%%%%%%%%%%%%%%%%
% %%%%%%%%%%%%%%%%%%%%%%%%%%%%%%%%%%%%%
% %%%%%%%%%%%%   SECTION   %%%%%%%%%%%%
% %%%%%%%%%%%%%%%%%%%%%%%%%%%%%%%%%%%%%
% %%%%%%%%%%%%%%%%%%%%%%%%%%%%%%%%%%%%%
% \section*{Journals}
% \begin{list}{}
% 	{\leftmargin=2em \itemindent=-2em}
% 	\item 
% 		\textsc{Alex the student, Adam, Eve}: ``Challenges in Online Learning Platforms'', EURASTUD Journal on European Student Teaching,  2008, P. 1\hbox{-}10.
% 	%\item
% \end{list}
% \smallskip

% %%%%%%%%%%%%%%%%%%%%%%%%%%%%%%%%%%%%%
% %%%%%%%%%%%%%%%%%%%%%%%%%%%%%%%%%%%%%
% %%%%%%%%%%%%   SECTION   %%%%%%%%%%%%
% %%%%%%%%%%%%%%%%%%%%%%%%%%%%%%%%%%%%%
% %%%%%%%%%%%%%%%%%%%%%%%%%%%%%%%%%%%%%
% \section*{International Conferences}
% \begin{list}{}
% 	{\leftmargin=2em \itemindent=-2em}
% 	\item 
% 		\textsc{Alex the student}: ``Practical Online Learning Tools'', in 21st International Student Teaching Conference (ISTC), 2012, P. 1\hbox{-}7.
% 	%\item
% \end{list}
% \smallskip

% %%%%%%%%%%%%%%%%%%%%%%%%%%%%%%%%%%%%%
% %%%%%%%%%%%%%%%%%%%%%%%%%%%%%%%%%%%%%
% %%%%%%%%%%%%   SECTION   %%%%%%%%%%%%
% %%%%%%%%%%%%%%%%%%%%%%%%%%%%%%%%%%%%%
% %%%%%%%%%%%%%%%%%%%%%%%%%%%%%%%%%%%%%
% \section*{Workshops}
% \begin{list}{}
% 	{\leftmargin=2em \itemindent=-2em}
% 	\item 
% 	\textsc{Alex the student, Eve} : ``Online Learning Monitoring and Debugging through Measurement Visualization'', in the 4th IEEE International Workshop on Hot Topics in Online Learning (IEEE HotOnLe'12), 2012, P. 1\hbox{-}6.
% 	%\item
% \end{list}
% \smallskip


% %%%%%%%%%%%%%%%%%%%%%%%%%%%%%%%%%%%%%
% %%%%%%%%%%%%%%%%%%%%%%%%%%%%%%%%%%%%%
% %%%%%%%%%%%%   SECTION   %%%%%%%%%%%%
% %%%%%%%%%%%%%%%%%%%%%%%%%%%%%%%%%%%%%
% %%%%%%%%%%%%%%%%%%%%%%%%%%%%%%%%%%%%%
% %\section*{Posters and Demos}
% %\begin{list}{}
% %	{\leftmargin=2em \itemindent=-2em}
% %	\item 
% %	\textsc{Alex the student, Adam}: ``OnTeMoS: Online Teacher Monitoring System'', in the 3rd International Workshop on Online Learning Approaches, Experimental Evaluation and Characterization, New York, NY, USA, 20013, P. 101\hbox{-}102.
% 	%\item
% %\end{list}

%EIDESSTATTLICHE ERKLAERUNG
%SELBSTAENDIGKEITSERKLAERUNG
\newcommand{\theauthor}{Elham Aref}
\pagestyle{empty}

\vspace*{\stretch{1}}
\setlength{\parindent}{0in}
%Ich versichere an Eides statt, dass ich diese Diplomarbeit selbst\"andig verfasst und nur die angegebenen Quellen und Hilfsmittel verwendet habe.
%
%Alle w\"ortlich oder inhaltlich \"ubernommenen Stellen habe ich als solche gekennzeichnet.
%
%Die vorliegende Arbeit wurde weder in der vorliegenden noch einer modifizierten Fassung einer dritten, in- oder ausl\"andischen Fakult\"at als Pr\"ufungsleistung, oder zum Erlangen eines akademischen Grades vorgelegt.

I hereby confirm this thesis been written on my own and only used the aids indicated.
The current work was not submitted to a third, domestic or foreign faculty as an examination, neither in the present nor in a modified version, or to obtain an academic degree.

\vspace{1cm}

\begin{flushright}
	\begin{tabular}{p{2cm}p{6cm}}
		\hline
		2023.2.14  & \theauthor
		\smallskip
	\end{tabular}
\end{flushright}


% %INSERT EMPTY PAGE
% \blankpage

%TOC PAGE
\tableofcontents

%MAIN CONTENT
\mainmatter \pagestyle{fancy}
%\pagestyle{headings}

\chapter{Introduction}
\label{chap:introduction}

General Introduction


\section{Problem Statement}
\label{sec:intro:probstatement}

\section{Contributions}
\label{sec:intro:contrib}


\begin{itemize}
\item {\bf Measurements}
\item{\bf Evaluation Tools}

\end{itemize}

%%%%%%%%%%%%%%%%%%%%%%%%%%%%%%%%%%%%%%%%%%%%%%%%%%%%%%%%%%%%%%%%%%%%%%%%%%
%%%%%%%%%%%%   CAPTER 2   %%%%%%%%%%%%%%%%%%%%%%%%%%%%%%%%%%%%%%%%%%%%%%%%
%%%%%%%%%%%%%%%%%%%%%%%%%%%%%%%%%%%%%%%%%%%%%%%%%%%%%%%%%%%%%%%%%%%%%%%%%%
\chapter{Related Work}
\label{chap:related_work}

%intro
%
To this end, we start with a short overview and then present the state-of-the-art of current approaches. Next, we discuss the current work.
%
Finally, we conclude with a summary and open problems.


%%%%%%%%%%%%%%%%%%%%%%%%%%%%%%%%%%%%%
%%%%%%%%%%%%%%%%%%%%%%%%%%%%%%%%%%%%%
%%%%%%%%%%%%   SECTION   %%%%%%%%%%%%
%%%%%%%%%%%%%%%%%%%%%%%%%%%%%%%%%%%%%
%%%%%%%%%%%%%%%%%%%%%%%%%%%%%%%%%%%%%
\section{A Short Overview}
\label{s:basics}

%\begin{figure}
%	\center
%	\includegraphics[width=\columnwidth]{figures/overview}
%	\caption{Overview.}
%	\label{fig:overview}
%\end{figure}


%%%%%%%%%%%%%%%%%%%%%%%%%%%%%%%%%%%%%
%%%%%%%%%%%%%%%%%%%%%%%%%%%%%%%%%%%%%
%%%%%%%%%%%%   SECTION   %%%%%%%%%%%%
%%%%%%%%%%%%%%%%%%%%%%%%%%%%%%%%%%%%%
%%%%%%%%%%%%%%%%%%%%%%%%%%%%%%%%%%%%%
\section{Control theory}
\label{s:single_control}

%
In this section, we discuss these approaches in their separate sections.
%

%%%%%%%%%%%%%%%%%%%%%%%%%
%%%%%   SUB-SECTION   %%%
%%%%%%%%%%%%%%%%%%%%%%%%%
%%%%%%%%%%%%%%%%%%%%%%%%%
\subsection{Old Control}
\label{ss:oc}


%summary table of old protocols
\begin{landscape}
\ctable[
	cap	= Summary of old Control (OC) Algorithms,
	caption = Summary of old Control (OC) Algorithms,
	label	= tab:oc_summary,
	pos  	= h
]{lccccc}{
%
\tnote[1]{Abbreviations: Obj: Objective, E: Energy, T: Topology, C: Capacity, Dist.: Distributed, Cent.:Centralized, synch: Synchronized, P: Prototype}
\tnote[2]{Old conditions are similar in space and time.}
\tnote[3]{Ideal comm: Overhearing any transmissions.}
\tnote[4]{New control considered.}
\tnote[5]{Angles between 5\textdegree to 120\textdegree are considered.}
%
} {
	\FL
	%
	\textbf{Protocol}			& \textbf{Obj.}\tmark[1]	& \textbf{Granularity}	& \textbf{Type of Control} & \textbf{PHY layer assumptions} & \textbf{Validation}\ML
	%
        PAR~\cite{Chen02-WNJ}			& E	& Per-packet	& Dist.	& Sym\tmark[2], ideal comm.\tmark[3],omni\tmark[4] & \no	\NN
        \hline	
	TER~\cite{154a}		& T	& Per-network	& Cent.	& Sym, ideal comm.,omni & \yes (Routing)	\NN
	MLP~\cite{aazhang88}		& T	& Per-cluster	& Cent.	& Sym, ideal comm.,omni & \yes (Routing)	\NN
        \hline	
	PCM~\cite{FuInfocom03}			& C	& Per-packet	& Dist.	& Sym, ideal comm., omni & \no	\NN
	POW~\cite{aiello03}	& C	& Per-packet	& Dist.	& Sym, ideal comm., omni & \no	\NN
	MID~\cite{callaway02}	& C	& Per packet	& Dist.	& Directional\tmark[5]	& \yes (MAC,WARP)	\NN
	DIR~\cite{bianchi00}			& C	& Per slot	& Cent.	& Directional	& \yes (MAC,Madwifi)	\NN
	SPE~\cite{crow97}			& C	& Per slot	& Cent.	& Directional	& \yes 	(Multi-radio)	\LL
}
\end{landscape}


In summary, it can be concluded that control does not bring energy savings.


%%%%%%%%%%%%%%%%%%%%%%%%%
%%%%%   SUB-SECTION   %%%
%%%%%%%%%%%%%%%%%%%%%%%%%
%%%%%%%%%%%%%%%%%%%%%%%%%
\subsection{New Control}
\label{ss:trc}

In this section, we summarize the algorithms in literature.
%
Table~\ref{tab:oc_summary} summarizes the approaches considered in this section.
%


%%%%%%%%%%%%%%%%%%%%%%%%%%%%%%%%%%%%%
%%%%%%%%%%%%%%%%%%%%%%%%%%%%%%%%%%%%%
%%%%%%%%%%%%   SECTION   %%%%%%%%%%%%
%%%%%%%%%%%%%%%%%%%%%%%%%%%%%%%%%%%%%
%%%%%%%%%%%%%%%%%%%%%%%%%%%%%%%%%%%%%
\section{Summary}

In this section, we gave an overview of the research on both theoretical and practical control.
%



\chapter{Measurement Tools}
\label{sec:mmeasurement}

\section{Rate Control API}
\label{sec:measurement:current}

\subsection{Capabilities}

Monitoring task commands  
start,stop,manual,auto ...
\label{sec:measurementg:validation}
\subsection{Setup}
How to take the traces
\subsection{Format of Monitoring Information}
\begin{sloppypar}
examples
% 
List of status reports
% 
format of trace rates with description.
% 
\begin{itemize}
\item \ phy0,timestamp,txs,mac
\item \ phy0,timestamp,rxs,mac
\item \ phy0,timestamp,stats,mac
\item \ phy0,timestamp,best-rates,mac
\end{itemize}
\end{sloppypar}



\section{Summary}


\chapter{Analyse and Optimization}
\label{chap:controller}


\section{Implementation of an Analysis Tool}
\label{sec:controller:design}

\section{Performance Evaluation}
\label{sec:controller:evaluation}


\section{Summary}
\label{sec:controller:conclusion}


\chapter{Summary and Outlook}
\label{chap:conclusion}

\section{Summary}


\section{Future Directions}



%BIBLIOGRAPHY
\addcontentsline{toc}{chapter}{Bibliography}
\bibliographystyle{IEEEtran}
\fancyhead[RE]{\normalfont Bibliography}
\fancyhead[LO]{\normalfont Bibliography}
\bibliography{bibliography/IEEEabrv,bibliography/sample.bib}


%ABBILDUNGSVERZEICHNIS
\addcontentsline{toc}{chapter}{List of Figures}
\listoffigures

%TABELLENVERZEICHNIS
\addcontentsline{toc}{chapter}{List of Tables}
\listoftables

%APPENDIX
%\appendix
%\chapter{Appendix}
\label{chap:appendix}


\backmatter
%\include{cv}

\end{document}
